\section{Discussion}\label{s:7}

\subsection{Why Adaptive Queue Control Outperforms Fixed Timing Under Asymmetric Disruptions}

The strongest performance gains of \texttt{dqn\_agent} occur in scenarios characterized by sharp spatial asymmetry or altered fleet composition, namely \texttt{80\_emergency\_multi\_ambulance} ($42.38\%$ queue reduction, $75.33\%$ accumulated-delay reduction), \texttt{64\_flooding\_dbp\_reduced} ($34.02\%$ queue reduction, $67.97\%$ accumulated-delay reduction), \texttt{30\_vip\_motorcade\_diversion} ($22.58\%$ queue reduction, $52.25\%$ accumulated-delay reduction), and \texttt{68\_tet\_return\_flood} ($24.47\%$ queue reduction, $60.14\%$ accumulated-delay reduction). Because the fixed-time program in \texttt{roundabout.net.xml} (\texttt{programID="1"}) allocates uniform $31\,\mathrm{s}$ intervals across all static phases without observing inbound densities, it continues serving empty or low-demand phases (such as $A_1=0$ in Scenario~30) while queues accumulate on heavily loaded approaches. By contrast, \texttt{MixedTrafficObservation} supplies per-approach PCU densities $\rho_i$, class-stratified queue vectors $\mathbf{q}_i$, and circulatory ring speeds $v_j$ every $5\,\mathrm{s}$, allowing the DQN policy to reallocate green combinations dynamically toward congested approaches while preserving circulatory flow.

In \texttt{49\_extreme\_overload\_all\_legs}, total demand reaches $23{,}000\,\mathrm{veh/h}$ across all five legs simultaneously. Under uniform global oversaturation, phase reallocation offers smaller marginal relief ($6.37\%$ queue reduction; $30.39\%$ accumulated waiting reduction) because granting right-of-way to one saturated approach immediately delays another. Similarly, in \texttt{77\_concert\_event\_district3}, where three approaches ($A_2, A_4, A_6$) carry balanced high loads, \texttt{fixed\_time} achieves slightly lower mean queue ($30.70$ vs.\ $31.60$ vehicles) and higher short-horizon arrivals ($4968.0$ vs.\ $4680.0\,\mathrm{veh/h}$), even though \texttt{dqn\_agent} still maintains lower accumulated waiting time ($10{,}983$ vs.\ $13{,}041\,\mathrm{s}$) and higher mean network speed ($9.73$ vs.\ $9.55\,\mathrm{m/s}$).

\subsection{Reconciling Completed-Trip Waiting Time with Network-Wide Queue and Delay}

A notable pattern in Table~\ref{tab:te_primary_results} is the apparent divergence between step-level telemetry and \texttt{tripinfo.xml}-derived averages in Scenarios~64 and~68: \texttt{dqn\_agent} cuts average queue length by $24.47\%\text{--}34.02\%$, reduces network-wide accumulated waiting time by $60.14\%\text{--}67.97\%$, and raises throughput by $12.86\%\text{--}19.38\%$, yet reports higher completed-trip waiting time ($4.82$ vs.\ $3.54\,\mathrm{s}$ in Scenario~64; $4.08$ vs.\ $2.19\,\mathrm{s}$ in Scenario~68) and time loss ($14.92$ vs.\ $11.22\,\mathrm{s}$; $13.33$ vs.\ $9.86\,\mathrm{s}$).

This discrepancy is explained by \emph{right-censoring (survivorship bias)} in \texttt{parse\_tripinfo} over the $100\,\mathrm{s}$ evaluation window. SUMO writes a \texttt{<tripinfo>} entry only when a vehicle finishes its route and exits the network before $t=100\,\mathrm{s}$. As shown in Fig.~\ref{fig:step-dashboard-64}, under \texttt{fixed\_time} in Scenario~64, queues rapidly climb to $76\text{--}80$ halted vehicles and network-wide accumulated waiting time reaches $19{,}768\,\mathrm{s}$, leaving many delayed vehicles trapped inside the network at $t=100\,\mathrm{s}$ (only $70$ vehicles, or $2520\,\mathrm{veh/h}$, exit). Consequently, \texttt{fixed\_time}'s \texttt{tripinfo.xml} averages reflect mainly unobstructed vehicles that entered early or traversed lightly loaded legs. By contrast, \texttt{dqn\_agent} actively discharges queued vehicles ($79$ completed trips, $+12.86\%$ in Scenario~64; $154$ vs.\ $129$, $+19.38\%$ in Scenario~68), bringing previously delayed vehicles into the completed-trip pool and raising the mean waiting time among arrivals while substantially lowering total network congestion.

\subsection{Role of Phase Switching and Reward Configuration}

In \texttt{dqn\_roundabout.yaml}, the flicker penalty weight is set to $w_{\text{flicker}} = 0.0$ while delay and queue weights are both $1.5$. Consequently, \texttt{dqn\_agent} receives no penalty for requesting frequent phase changes, exhibiting a mean switch rate of $0.58$ requests per $5\,\mathrm{s}$ step (up to $0.75$ in Scenario~49). Although the signal controller prevents unsafe switching before $t_{\min}=10\,\mathrm{s}$ ($0$ violations), frequent transitions trigger $3\,\mathrm{s}$ yellow clearance intervals, which introduce short deceleration--acceleration cycles. This explains why \texttt{dqn\_agent} achieves substantial reductions in instantaneous queues ($21.15\%$ mean per-scenario reduction) and accumulated waiting time ($51.52\%$ reduction in six-scenario mean total; $50.31\%$ mean per-scenario reduction) while six-scenario mean fuel ($+0.30\%$; $-0.29\%$ mean per-scenario improvement), $\mathrm{CO}_2$ ($+0.30\%$; $-0.30\%$ mean per-scenario improvement), and $\mathrm{NO}_x$ ($+1.67\%$; $-1.88\%$ mean per-scenario improvement) increase slightly relative to \texttt{fixed\_time}: stop-and-go waves induced by $3\,\mathrm{s}$ clearance transitions offset the emission benefits of shorter queues. Enabling a non-zero $w_{\text{flicker}}$ (as configured in \texttt{ppo\_roundabout.yaml} with $w_{\text{flicker}}=1.0$) represents a direct lever for trading marginal queue responsiveness for smoother phase stability and lower emissions.

\subsection{Operational Feasibility for Mixed-Traffic Deployment}

From an embedded-control perspective, the MLP policy requires $0.452\text{--}0.529\,\mathrm{ms}$ per inference step, well within the $5\,\mathrm{s}$ decision cycle and leaving ample headroom for roadside sensor aggregation and V2I/loop-detector preprocessing. Furthermore, the centralized state representation naturally accounts for motorcycle-dominated Vietnamese traffic (60\%--90\% two-wheelers across \texttt{TE} scenarios) via PCU weighting ($w_{\text{pcu}}=0.3$ for motorcycles vs.\ $2.0$ for buses and $2.5$ for trucks), preventing four-wheel queues from being masked by high two-wheeler counts.
